\documentclass[11pt]{article}
\usepackage[margin=1in]{geometry}
\usepackage{amsmath,amssymb,amsthm}
\usepackage{booktabs}
\usepackage{graphicx}
\usepackage{longtable}
\usepackage{array}
\usepackage{hyperref}
\usepackage{enumitem}
\usepackage{caption}
\usepackage{algorithm}
\usepackage{algorithmic}
\usepackage{mathtools}
\usepackage{xcolor}
% multirow removed (not available on this system)
\captionsetup{font=small}

\newtheorem{theorem}{Theorem}
\newtheorem{proposition}[theorem]{Proposition}
\newtheorem{lemma}[theorem]{Lemma}
\newtheorem{corollary}[theorem]{Corollary}
\newtheorem{definition}[theorem]{Definition}
\newtheorem{remark}[theorem]{Remark}
\newtheorem{example}[theorem]{Example}

\DeclareMathOperator{\TIP}{TIP}
\DeclareMathOperator{\supp}{supp}
\DeclareMathOperator{\Rips}{Rips}
\DeclareMathOperator{\diam}{diam}
\DeclareMathOperator{\Tr}{Tr}
\DeclareMathOperator*{\argmin}{arg\,min}

\title{CC-SPR: Geometry-Sensitive Sparse Persistent Representatives\\for Omics Data Analysis\\[6pt]
\large A Backend-Aware Extension of Harmonic Persistent Homology}
\author{}
\date{}

\begin{document}
\maketitle

\begin{abstract}
Persistent homology provides a multiscale summary of the topology of point-cloud data, but its standard output---a barcode or persistence diagram---discards the geometric identity of the cycles it detects.
Harmonic persistent homology recovers that identity by associating to each bar a representative cycle that minimizes an $\ell_2$ energy, creating a bridge from topological features back to individual samples and, via projection, to molecular features.
We introduce \textbf{CC-SPR} (Cycle-Consistent Sparse Persistent Representatives), an extension that modifies two layers of this pipeline.
First, we replace the dense $\ell_2$ representative with a sparse elastic-net objective that concentrates cycle support on a small number of edges, yielding representatives whose projected feature influence is more interpretable and more stable under bootstrap resampling.
Second, we formalize the sample-space metric as a \emph{backend} parameter, so that the same filtration--representation--projection pipeline can be instantiated with Euclidean distances, Ollivier--Ricci curvature-adjusted distances, diffusion distances, PHATE-like potential distances, or distance-to-measure (DTM) robust distances.

We evaluate CC-SPR on four cancer genomics datasets under a resource-bounded execution policy.
TCGA-LUAD serves as the anchor benchmark with multiseed paired statistics across 18 random splits; GSE161711 provides disease-aligned bulk CLL validation with full six-axis ablation; BRCA contributes supplementary support; and GSE165087 provides bounded single-cell heterogeneity validation after three progressively smaller configurations exhausted available cluster memory.
In the current implementation, Euclidean and Ollivier--Ricci backends are fully benchmarked; diffusion, PHATE-like, and DTM backends are mathematically developed as justified extensions whose empirical instantiation remains future work.
The main contribution is not a claim of universal predictive superiority---standard non-topological baselines remain strongest on several datasets---but rather a geometry-sensitive topological framework with explicit sparsity control, bootstrap-aggregated topological influence profiles, and an honest execution envelope that separates completed evidence from analytically motivated extensions.
\end{abstract}

% ============================================================
\section{Introduction}
\label{sec:introduction}
% ============================================================

\subsection{Omics data as a geometry-rich, noisy, high-dimensional setting}

Modern omics experiments---bulk RNA-seq, single-cell transcriptomics, methylation arrays, proteomics---routinely produce data matrices with $n$ samples in $10^1$--$10^5$ and $p$ features in $10^3$--$10^5$.
These matrices are not merely large: they encode geometry that reflects biological structure.
Tumor subtypes form clusters; differentiation trajectories trace manifold-like paths; regulatory circuits create loop-like patterns in expression space; and microenvironment heterogeneity introduces curvature and density variation.
At the same time, omics data is noisy---technical batch effects, dropout in single-cell protocols, measurement imprecision---and high-dimensional in ways that challenge standard Euclidean assumptions.

This combination of geometric richness and statistical noise makes omics a natural setting for topological data analysis (TDA), which provides multiscale structural summaries that are robust to continuous deformations of the data.
Persistent homology, the workhorse of TDA, has been applied to gene expression landscapes \cite{rizvi2017}, protein structure \cite{xia2014}, tumor microenvironments \cite{vipond2021}, single-cell trajectories \cite{moon2019phate}, neural population coding \cite{giusti2015}, and viral evolution \cite{chan2013topology}.

\subsection{Why persistent homology matters for omics}

The appeal of persistent homology in omics is that it captures collective structure---loops, voids, connected components---that cannot be detected by feature-wise statistics or pairwise correlations alone.
A cluster of samples sharing a gene-expression phenotype is a $H_0$ feature; a loop in sample space that corresponds to a regulatory feedback circuit or a cell-cycle trajectory is an $H_1$ feature.
The persistence of such features across filtration scales provides a natural measure of their robustness: short-lived features are likely noise, while long-lived features may reflect genuine biological signal.

However, the standard persistent-homology pipeline has a well-known interpretability gap: while a barcode tells us \emph{that} a topological feature exists and \emph{how long} it persists, it does not tell us \emph{which samples} participate in the corresponding cycle, nor \emph{which molecular features} drive it.
This gap limits the biological utility of TDA in omics, where the end goal is typically to identify genes, pathways, or modules associated with a structural signal.

\subsection{Why harmonic persistent homology is important}

Harmonic persistent homology \cite{harmonicPH} addresses this interpretability gap by computing a representative cycle for each persistent class that minimizes an $\ell_2$ (harmonic) energy subject to homological constraints.
The representative can then be projected back to feature space, creating a bridge from topology to biology: a persistent loop in sample space can be traced to a set of genes whose coordinated variation defines that loop.
This is a significant conceptual advance because it transforms persistent homology from a shape descriptor into a feature-attribution tool.

\subsection{Limitations of the harmonic approach}

Two practical limitations motivate the present work.

\paragraph{Limitation 1: dense representatives blur interpretation.}
The harmonic $\ell_2$ objective spreads weight broadly across the edge support of a cycle class.
In omics, where the goal is often to identify a compact set of genes or pathways associated with a topological signal, dense support complicates downstream interpretation.
A cycle that distributes weight uniformly over 200 edges is harder to connect to a biological hypothesis than one that concentrates weight on 15 edges.
This is not a theoretical nicety: in practice, dense representatives project to diffuse feature-importance vectors that implicate many genes at similar low weights, making pathway analysis and biological interpretation ambiguous.

\paragraph{Limitation 2: Euclidean geometry may be inappropriate.}
The default sample-space metric in most persistent-homology pipelines is Euclidean distance in the ambient feature space or a PCA projection thereof.
This assumption is often too rigid for biological data with nonuniform density, local bridges between subtypes, curvature-rich manifold structure, or progression-like developmental trajectories.
Different geometry backends can produce different filtrations, different barcodes, and different representatives---a sensitivity that should be modeled explicitly rather than ignored.
When a Euclidean neighborhood graph contains shortcut edges between biologically distinct subpopulations, the resulting persistent features may be artifacts of metric rigidity rather than reflections of genuine topology.

\subsection{Why sparse, geometry-aware representatives matter biologically}

Sparse representatives matter in omics because they directly address the question biologists care about: \emph{which genes are implicated by a topological signal?}
A sparse cycle concentrates weight on a small number of edges, which project to a compact gene set.
This compact gene set is more amenable to pathway analysis, more interpretable as a regulatory module, and more testable as a biomarker panel than a diffuse weight distribution.

Geometry-aware representatives matter because the choice of sample-space metric determines which topological features are detected and which are suppressed.
Curvature-adjusted metrics can sharpen boundaries between subtypes; diffusion metrics can denoise local perturbations; robust metrics can suppress outlier-driven artifacts.
Making geometry a selectable backend---rather than a fixed assumption---enables systematic study of how metric choice affects biological interpretation.

\subsection{Contributions}

CC-SPR addresses both limitations within the harmonic persistent homology scaffold:
\begin{enumerate}[leftmargin=1.5em]
\item \textbf{Sparse elastic representatives.} We replace the dense $\ell_2$ objective with a convex elastic-net objective that encourages compact edge support while maintaining numerical stability (Section~\ref{sec:sparse}).
\item \textbf{Backend-aware geometry.} We formalize the sample-space metric as a selectable backend, developing Euclidean, Ollivier--Ricci, diffusion, PHATE-like, and DTM geometry families within a unified filtration-perturbation framework (Sections~\ref{sec:backends}--\ref{sec:dtm}).
\item \textbf{Topological influence profiles (TIP).} We introduce a bootstrap-aggregated feature-stability diagnostic that converts sparse representative output into a per-feature selection frequency, enabling direct comparison of support concentration across geometries (Section~\ref{sec:tip}).
\item \textbf{Resource-bounded evaluation.} We present a four-dataset evaluation under an explicit execution envelope, with paired multiseed statistics on the anchor benchmark and honest reporting of resource failures on single-cell data (Sections~\ref{sec:experimental}--\ref{sec:results}).
\end{enumerate}

\paragraph{What CC-SPR is not.}
CC-SPR is not a claim that Ricci curvature universally improves predictive accuracy, nor that topology should replace conventional classifiers.
It is a framework that makes topology-derived features more interpretable and more geometry-aware.
The value is diagnostic and methodological: sparse representatives with explicit geometry control yield a richer analysis object than dense Euclidean representatives alone.
This paper is an analytical methods study about structure, stability, interpretability, and topology-aware feature attribution---not a benchmark-chasing leaderboard paper.

% ============================================================
\section{Background: Persistent Homology and Harmonic Representatives}
\label{sec:background}
% ============================================================

\subsection{Simplicial homology and persistence}

Let $X = \{x_1, \ldots, x_n\} \subset \mathbb{R}^p$ be a point cloud equipped with a metric $d$.
The \emph{Vietoris--Rips complex} at scale $\varepsilon$ is
\[
\Rips_\varepsilon(X, d) = \bigl\{ \sigma \subseteq X : \diam_d(\sigma) \le \varepsilon \bigr\}.
\]
As $\varepsilon$ increases from $0$ to $\infty$, the nested family $\{\Rips_\varepsilon\}_{\varepsilon \ge 0}$ forms a filtration.
Applying simplicial homology with field coefficients $\mathbb{F}$ to this filtration yields a persistence module.
By the structure theorem for graded modules over a PID, this module decomposes into a direct sum of interval modules, each corresponding to an interval $[b_i, d_i)$ in the barcode \cite{zomorodian2005}.

\begin{definition}[Persistence barcode]\label{def:barcode}
The \emph{barcode} of a filtered simplicial complex is the multiset of intervals $\{[b_i, d_i)\}_{i \in I}$ encoding the birth and death scales of homology classes.
The \emph{persistence} (or \emph{lifetime}) of the $i$-th feature is $\pi_i = d_i - b_i$, and its \emph{midlife} is $m_i = (b_i + d_i)/2$.
\end{definition}

\begin{definition}[Chain groups and boundary maps]\label{def:chains}
Let $K$ be a simplicial complex.
The \emph{$k$-chain group} $C_k(K; \mathbb{F})$ is the $\mathbb{F}$-vector space with basis the $k$-simplices of $K$.
The \emph{boundary operator} $\partial_k: C_k \to C_{k-1}$ is the linear map defined on each $k$-simplex $\sigma = [v_0, \ldots, v_k]$ by
\[
\partial_k(\sigma) = \sum_{i=0}^k (-1)^i [v_0, \ldots, \hat{v}_i, \ldots, v_k],
\]
where $\hat{v}_i$ denotes omission of vertex $v_i$.
The fundamental identity $\partial_{k-1} \circ \partial_k = 0$ ensures that $\operatorname{im}(\partial_{k+1}) \subseteq \ker(\partial_k)$, and the $k$-th homology group is $H_k(K; \mathbb{F}) = \ker(\partial_k) / \operatorname{im}(\partial_{k+1})$.
\end{definition}

\subsection{Harmonic representatives}

For a selected $H_1$ bar $[b, d)$, let $z_0 \in C_1(\Rips_\varepsilon; \mathbb{R})$ be an initial cycle representative at scale $\varepsilon \in [b, d)$, and let $\partial_2 : C_2 \to C_1$ denote the boundary operator mapping 2-chains (triangles) to 1-chains (edges).
The full set of cycles homologous to $z_0$ forms the \emph{affine homology class}
\begin{equation}\label{eq:affine}
[z_0] + \operatorname{im}(\partial_2) = \{z_0 + \partial_2 y : y \in C_2\}.
\end{equation}

The \emph{harmonic representative} minimizes the $\ell_2$ energy within this affine class:
\begin{equation}\label{eq:harmonic}
\hat{z}_{\mathrm{harm}} = \argmin_{y \in C_2} \| z_0 + \partial_2 y \|_2^2.
\end{equation}
This is a standard least-squares problem whose solution is the orthogonal projection of $z_0$ onto $\ker(\partial_1) \cap (\operatorname{im} \partial_2)^\perp$, coinciding with the harmonic representative in the Hodge-theoretic sense when the chain complex is equipped with the standard inner product \cite{harmonicPH, friedman1998}.

\begin{remark}[Dense support problem]\label{rem:dense}
The solution to \eqref{eq:harmonic} typically has $\|\hat{z}_{\mathrm{harm}}\|_0 = O(|E|)$ where $|E|$ is the number of edges at scale $\varepsilon$.
In omics applications where $n$ ranges from $10^2$ to $10^4$, this can mean hundreds or thousands of nonzero edge weights, diluting the biological signal when projected to feature space.
The harmonic representative finds the smoothest cycle, not the most parsimonious one.
\end{remark}

\subsection{Conceptual bridge from harmonic PH to CC-SPR}

CC-SPR preserves the four-step harmonic persistent homology scaffold:
\begin{enumerate}[leftmargin=1.5em]
\item construct a sample-space metric,
\item compute a filtration and persistent homology,
\item choose a representative cycle for a selected homology class,
\item project representative structure back to features for downstream learning or interpretation.
\end{enumerate}
The extension is targeted rather than wholesale: Step~3 is replaced by a sparse elastic representative, and Step~1 is lifted from a fixed Euclidean choice to a backend family.
This makes CC-SPR a strict extension of the harmonic framework, not a disconnected competing method.
Every harmonic PH analysis is a special case of CC-SPR with $\lambda_1 = 0$ and a Euclidean backend.

\subsection{Comparison of harmonic PH and CC-SPR}

Table~\ref{tab:hphcompare} provides a systematic comparison across eight axes relevant to omics applications.

\begin{table}[h]
\centering
\small
\begin{tabular}{p{2.5cm}p{3.0cm}p{3.0cm}}
\toprule
\textbf{Axis} & \textbf{Harmonic PH} & \textbf{CC-SPR} \\
\midrule
Metric assumption & Fixed Euclidean & Selectable backend family \\
Representative criterion & Dense $\ell_2$ (harmonic) & Sparse elastic net \\
Sparsity & No explicit control & $\lambda_1$-controlled \\
Interpretability & Diffuse feature weights & Concentrated feature weights \\
Stability diagnostics & Not built-in & TIP profiles (bootstrap) \\
Computational cost & Moderate & Moderate to high (backend-dependent) \\
Backend flexibility & None (single metric) & Euclidean, Ricci, diffusion, PHATE, DTM \\
Biological alignment & Medium (diluted signal) & High (compact gene sets) \\
Implementation status & Reference framework & Eu/Ricci implemented; others mathematical \\
\bottomrule
\end{tabular}
\caption{Systematic comparison of harmonic persistent homology and CC-SPR across axes relevant to omics analysis.
CC-SPR extends harmonic PH by adding explicit sparsity control, stability diagnostics, and geometry-backend awareness while preserving the same filtration--representation--projection scaffold.}
\label{tab:hphcompare}
\end{table}

\paragraph{When is the harmonic representative sufficient?}
The harmonic representative is sufficient when (i) the cycle support is naturally compact (e.g., small $H_1$ bars with few participating edges), (ii) the downstream task does not require feature-level attribution (e.g., barcode-only analyses), or (iii) the dataset is small enough that dense support does not dilute biological signal.

\paragraph{When does CC-SPR add value?}
CC-SPR adds value in three scenarios:
(i)~when the analysis goal requires identifying a compact gene set associated with a topological feature, since sparse representatives concentrate weight on fewer edges and thus fewer genes;
(ii)~when geometry sensitivity is itself a scientific question, since the backend-aware framework enables controlled comparison of how metric choice affects topological features;
(iii)~when bootstrap stability of feature selection matters, since TIP profiles provide a principled stability diagnostic unavailable in the standard harmonic PH pipeline.

% ============================================================
\section{Mathematical Formulation}
\label{sec:math}
% ============================================================

\subsection{Notation and setup}

Let $X \in \mathbb{R}^{n \times p}$ denote an omics data matrix with $n$ samples (rows) and $p$ features (columns).
In bulk RNA-seq, rows are patients and columns are genes; in single-cell RNA-seq, rows are cells and columns are genes or principal components.
A \emph{geometry backend} $\mathcal{G}$ with hyperparameters $\theta_{\mathcal{G}}$ induces a sample-space distance matrix
\[
D^{(\mathcal{G})} \in \mathbb{R}_{\ge 0}^{n \times n}, \qquad D^{(\mathcal{G})}_{ii}=0, \qquad D^{(\mathcal{G})}_{ij} = D^{(\mathcal{G})}_{ji}.
\]

\subsection{Weighted graph construction and metric deformation}

From $D^{(\mathcal{G})}$, we construct a weighted $k$-nearest-neighbor graph $G = (V, E, w)$ where $V = \{1, \ldots, n\}$, $E$ contains edges between $k$-nearest neighbors, and $w_{ij} = D^{(\mathcal{G})}_{ij}$.
For geometry backends that apply a flow or deformation (e.g., Ricci flow), the edge weights are iteratively updated:
\[
w_{ij}^{(t+1)} = f(w_{ij}^{(t)}, \kappa^{(t)}(i,j), \eta),
\]
where $f$ is the update rule, $\kappa^{(t)}$ is the curvature at iteration $t$, and $\eta$ is a step size.
The final distance matrix $D^{(\mathcal{G})}$ is obtained as shortest-path distances on the deformed graph.

\subsection{Vietoris--Rips filtration and persistence}

The Vietoris--Rips filtration on $D^{(\mathcal{G})}$ produces a nested sequence of simplicial complexes:
\[
\emptyset = \Rips_0 \subseteq \Rips_{\varepsilon_1} \subseteq \cdots \subseteq \Rips_{\varepsilon_N} = \Delta^{n-1},
\]
where $\Delta^{n-1}$ is the full simplex.
Applying homology with field coefficients yields a persistence module that decomposes into bars $\{[b_i, d_i)\}_{i \in I}$.
For CC-SPR, we focus on $H_1$ bars because 1-cycles in sample space can be projected to feature space via the data matrix, creating a direct bridge from topological structure to gene-level influence.

\subsection{Representative classes and the affine constraint}

For a selected $H_1$ bar $[b, d)$, any representative cycle lies in the affine homology class
\[
z = z_0 + \partial_2 y, \qquad y \in C_2,
\]
where $z_0$ is an initial representative and $\partial_2$ is the boundary operator from 2-chains to 1-chains.
The vector $z \in \mathbb{R}^{|E|}$ assigns a weight to each edge; $z_e \ne 0$ means edge $e$ participates in the cycle with weight $z_e$.

\subsection{Harmonic $\ell_2$ objective}

The harmonic representative minimizes
\begin{equation}\label{eq:harm2}
\hat{y}_{\mathrm{harm}} = \argmin_{y \in C_2} \| z_0 + \partial_2 y \|_2^2.
\end{equation}
This is a convex quadratic with a closed-form solution via the normal equations $\partial_2^\top \partial_2 \hat{y} = -\partial_2^\top z_0$.

\subsection{Sparse elastic objective}
\label{sec:sparse}

CC-SPR replaces \eqref{eq:harm2} with
\begin{equation}\label{eq:ccspr}
\hat{y}_{\mathrm{spr}} = \argmin_{y \in C_2} \;\underbrace{\| z_0 + \partial_2 y \|_2^2}_{\text{fidelity}} + \underbrace{\lambda_1 \| z_0 + \partial_2 y \|_1}_{\text{sparsity}} + \underbrace{\lambda_2 \| y \|_2^2}_{\text{stability}},
\end{equation}
with $\lambda_1, \lambda_2 \ge 0$.
The sparse representative is $\hat{z}_{\mathrm{spr}} = z_0 + \partial_2 \hat{y}_{\mathrm{spr}}$.

\begin{proposition}[Convexity]\label{prop:convex}
The objective in \eqref{eq:ccspr} is convex in $y$.
When $\lambda_2 > 0$, it is strongly convex, guaranteeing a unique minimizer.
\end{proposition}

\begin{proof}
Each term is a convex function of $y$:
(i) $f_1(y) = \|z_0 + \partial_2 y\|_2^2$ is convex as a composition of a convex quadratic with the affine map $y \mapsto z_0 + \partial_2 y$;
(ii) $f_2(y) = \lambda_1 \|z_0 + \partial_2 y\|_1$ is convex because the $\ell_1$ norm is convex and composition with affine maps preserves convexity;
(iii) $f_3(y) = \lambda_2 \|y\|_2^2$ is convex (strongly convex when $\lambda_2 > 0$).
A nonnegative sum of convex functions is convex, and strong convexity of any summand implies strong convexity of the sum.
\end{proof}

\begin{remark}[Connection to elastic net regression]\label{rem:elasticnet}
The structure of \eqref{eq:ccspr} mirrors the elastic net \cite{zou2005elasticnet}: $\lambda_1$ plays the role of the LASSO penalty promoting sparsity in the edge weights, while $\lambda_2$ provides the ridge-like stability that prevents ill-conditioning when $\partial_2$ is rank-deficient.
The key difference is that the ``design matrix'' here is the boundary operator $\partial_2$ rather than a data matrix, and the ``response'' is the initial cycle $z_0$.
\end{remark}

\begin{proposition}[Sparsity guarantee]\label{prop:sparsity}
For any $\lambda_1 > 0$, the solution $\hat{z}_{\mathrm{spr}}$ satisfies
\[
\|\hat{z}_{\mathrm{spr}}\|_0 \le \|\hat{z}_{\mathrm{harm}}\|_0,
\]
with equality only when the harmonic representative already has minimal $\ell_1$ norm in its homology class.
\end{proposition}

\begin{proof}[Sketch]
The $\ell_1$ penalty acts as a soft-thresholding operator on the edge weights of the representative.
Any edge weight $|z_e| < \tau(\lambda_1)$ for an appropriate threshold $\tau$ is driven to zero.
Since the harmonic representative has no such thresholding, its support is at least as large.
Equality requires that every nonzero weight in $\hat{z}_{\mathrm{harm}}$ exceeds $\tau(\lambda_1)$, which occurs only when the harmonic solution already has minimal $\ell_1$ norm---a non-generic condition.
\end{proof}

\subsection{Topological influence profile (TIP)}
\label{sec:tip}

\begin{definition}[Topological Influence Profile]\label{def:tip}
For bootstrap replicates $b = 1, \ldots, B$ and a top-$K$ feature selection threshold, let $S_b \subseteq \{1, \ldots, p\}$ denote the set of features with the $K$ largest absolute projected weights from the sparse representative in replicate $b$.
The \emph{topological influence profile} is
\[
\TIP(j) = \frac{1}{B} \sum_{b=1}^B \mathbf{1}\{j \in S_b\}, \qquad j = 1, \ldots, p.
\]
\end{definition}

\begin{remark}[TIP as a stability diagnostic]\label{rem:tipstability}
Concentrated TIP mass (a few features with $\TIP(j) \approx 1$) indicates that the topological signal consistently localizes onto a compact feature subset under resampling.
Diffuse TIP mass indicates unstable or distributed support.
Because the same TIP computation applies to any geometry backend, it provides a natural diagnostic for comparing how different backends affect the \emph{interpretability} of topological features, independent of downstream predictive performance.
\end{remark}

\subsection{Representative support diagnostics}
\label{sec:diagnostics_math}

Beyond TIP, the quality of a sparse representative can be characterized by several support diagnostics.

\begin{definition}[Support size]\label{def:support}
For a representative $\hat{z} \in \mathbb{R}^{|E|}$, the \emph{support size} is $\|\hat{z}\|_0 = |\{e \in E : \hat{z}_e \ne 0\}|$.
After projection to feature space, the \emph{projected support size} is $|\{j \in \{1,\ldots,p\} : |w_j| > 0\}|$ where $w \in \mathbb{R}^p$ is the projected feature-importance vector.
\end{definition}

\begin{definition}[Representative entropy and Gini coefficient]\label{def:entropy}
Let $\hat{p}_e = |\hat{z}_e| / \sum_{e'} |\hat{z}_{e'}|$ be the normalized absolute edge weights.
The \emph{representative entropy} is
\[
H(\hat{z}) = -\sum_{e : \hat{z}_e \ne 0} \hat{p}_e \log \hat{p}_e.
\]
Low entropy indicates concentrated support; high entropy indicates diffuse support.
The \emph{Gini coefficient} of the weight distribution provides a complementary inequality measure:
\[
G(\hat{z}) = \frac{\sum_{i=1}^m \sum_{j=1}^m ||\hat{z}_{e_i}| - |\hat{z}_{e_j}||}{2m \sum_{i=1}^m |\hat{z}_{e_i}|},
\]
where $m = \|\hat{z}\|_0$ is the support size.
High Gini indicates that weight is concentrated on a few edges; low Gini indicates uniform distribution.
\end{definition}

\begin{remark}[Support concentration and biological interpretability]\label{rem:concentration}
In omics, a representative with high Gini coefficient and low entropy concentrates its topological signal on a small number of sample-pair relationships, which project to a compact gene set.
This concentration is what makes the representative biologically actionable: it identifies a gene module rather than a diffuse background signal.
The $\ell_1$ penalty in \eqref{eq:ccspr} directly promotes this concentration by driving small weights to zero.
\end{remark}

\begin{remark}[Per-bootstrap entropy and support-size diagnostics]\label{rem:diagnostics_note}
In the current repository, per-bootstrap support vectors were not cached at sufficient granularity to compute post hoc entropy and Gini coefficients.
These diagnostics would require rerunning topology loops.
TIP, $H_1$ lifetime, and $H_1$ prominence are available from cached outputs and serve as the primary support diagnostics in this paper.
Full entropy/Gini analysis remains future work.
\end{remark}

\subsection{Stability under bounded metric perturbation}

\begin{proposition}[Filtration stability under backend perturbation]\label{prop:stability}
Let $D^{(a)}$ and $D^{(b)}$ be distance matrices from two backends.
Define $\varepsilon_\infty = \|D^{(a)} - D^{(b)}\|_\infty$.
Then the Rips filtrations satisfy an $\varepsilon_\infty$-interleaving:
\[
\Rips_r(X, D^{(a)}) \subseteq \Rips_{r + \varepsilon_\infty}(X, D^{(b)}) \subseteq \Rips_{r + 2\varepsilon_\infty}(X, D^{(a)})
\]
for all $r \ge 0$.
Consequently, the bottleneck distance between the corresponding persistence diagrams satisfies
\[
d_B\bigl(\mathrm{Dgm}(D^{(a)}),\, \mathrm{Dgm}(D^{(b)})\bigr) \le \varepsilon_\infty.
\]
\end{proposition}

\begin{proof}
The first inclusion follows because if $\diam_{D^{(a)}}(\sigma) \le r$, then $\diam_{D^{(b)}}(\sigma) \le r + \varepsilon_\infty$ by the definition of $\varepsilon_\infty = \max_{i,j} |D^{(a)}_{ij} - D^{(b)}_{ij}|$.
The second inclusion follows symmetrically.
The bottleneck bound is the standard algebraic stability theorem for Rips filtrations applied to $\varepsilon_\infty$-interleaved persistence modules \cite{chazal2014persistence}.
\end{proof}

\begin{remark}[Practical implication]\label{rem:perturbation}
This proposition means that backend choice can be studied as a \emph{controlled perturbation} of the topological output.
When $\varepsilon_\infty$ is small relative to persistence values, the barcode is approximately preserved; when it is large, the backend is genuinely changing the topological signal.
The paired Ricci-vs-Euclidean statistics reported in Section~\ref{sec:results} can be understood through this lens: the $\varepsilon_\infty$ between Euclidean and Ricci distance matrices on LUAD is large enough to shift some bars, producing the observed $+0.046$ mean delta in downstream classifier performance.
\end{remark}

\begin{remark}[Bridge edges, bottleneck sensitivity, and backend sharpening]\label{rem:bridge}
Backends that preferentially upweight bottleneck or bridge-like edges---as Ricci flow does for negatively curved edges---can sharpen the persistent signal by increasing the filtration parameter at which spurious shortcuts create topological noise.
This is the mechanism by which curvature-aware backends can improve representative quality: by delaying the merger of distinct topological components, they preserve genuine cycles at the expense of noise-induced ones.
The effect is data-dependent and non-monotone, as demonstrated by the divergent LUAD/GSE161711 ablation results.
\end{remark}

\subsection{Algorithm}

\begin{algorithm}[h]
\caption{CC-SPR Pipeline}
\label{alg:ccspr}
\begin{algorithmic}[1]
\REQUIRE Data matrix $X \in \mathbb{R}^{n \times p}$, geometry backend $\mathcal{G}$, parameters $\lambda_1, \lambda_2, k, K$, bootstraps $B$
\ENSURE Topological influence profile $\TIP \in [0,1]^p$
\STATE Compute distance matrix $D^{(\mathcal{G})} \gets \texttt{Backend}(X, \mathcal{G})$
\STATE Compute Vietoris--Rips filtration and persistent homology on $D^{(\mathcal{G})}$
\STATE Select the most persistent $H_1$ bar $[b, d)$
\STATE Extract initial cycle $z_0$ and boundary operator $\partial_2$ at scale $\varepsilon = (b + d)/2$
\FOR{$b = 1, \ldots, B$}
    \STATE Draw bootstrap resample $X^{(b)}$ of size $n$
    \STATE Recompute $D^{(\mathcal{G})(b)}$, filtration, and persistent homology
    \STATE Solve \eqref{eq:ccspr} to obtain $\hat{z}_{\mathrm{spr}}^{(b)}$
    \STATE Project $\hat{z}_{\mathrm{spr}}^{(b)}$ to feature space: $w^{(b)} \in \mathbb{R}^p$
    \STATE Select top-$K$ features: $S_b \gets \{j : |w^{(b)}_j| \text{ among top-}K\}$
\ENDFOR
\STATE Compute $\TIP(j) \gets \frac{1}{B} \sum_{b=1}^B \mathbf{1}\{j \in S_b\}$ for each feature $j$
\RETURN $\TIP$
\end{algorithmic}
\end{algorithm}

\subsection{Computational complexity}

Let $n$ denote the number of samples.
The dominant costs are:
\begin{itemize}[leftmargin=1.5em]
\item \textbf{Distance matrix}: $O(n^2 p)$ for Euclidean; $O(n^2 p + T \cdot |E_{\text{graph}}|)$ for Ricci flow with $T$ iterations.
\item \textbf{Persistent homology}: $O(m^3)$ worst case where $m$ is the number of simplices; practical implementations are much faster \cite{gudhi2014}.
\item \textbf{Sparse optimization}: $O(|E|^2 \cdot |F|)$ per bootstrap, where $|E|$ and $|F|$ are the number of edges and triangles at the chosen scale.
\item \textbf{Total}: $O(B \cdot (n^2 p + m^3))$, dominated by the bootstrap loop.
\end{itemize}

% ============================================================
\section{Geometry Backends}
\label{sec:backends}
% ============================================================

A central design principle of CC-SPR is that the sample-space metric is a \emph{selectable parameter}, not a fixed assumption.
We formalize this through the filtration-perturbation framework established in Proposition~\ref{prop:stability}.

\begin{definition}[Geometry backend]\label{def:backend}
A \emph{geometry backend} is a pair $(\mathcal{G}, \theta_{\mathcal{G}})$ consisting of a metric construction algorithm $\mathcal{G}$ and hyperparameters $\theta_{\mathcal{G}}$ that maps a data matrix $X \in \mathbb{R}^{n \times p}$ to a symmetric distance matrix $D^{(\mathcal{G})} \in \mathbb{R}_{\ge 0}^{n \times n}$ with $D^{(\mathcal{G})}_{ii} = 0$.
\end{definition}

\subsection{Euclidean geometry (baseline)}

The baseline backend uses the standard $\ell_2$ distance in PCA space:
\[
D^{(\text{eu})}_{ij} = \|x_i - x_j\|_2.
\]
Its advantages are simplicity, interpretability, and computational efficiency.
Its limitation is rigidity: it treats all directions equally and does not adapt to local density or curvature.
\emph{Implementation status: fully implemented and benchmarked.}

\subsection{Ricci curvature, bridge edges, and bottleneck sharpening}
\label{sec:ricci}

\begin{definition}[Ollivier--Ricci curvature on graphs \cite{ollivier2007,ollivier2009}]\label{def:ricci}
Let $G = (V, E, w)$ be a weighted graph.
For each vertex $v$, define a probability measure $\mu_v$ supported on its neighbors (e.g., uniform on the $k$-nearest neighbors).
The \emph{Ollivier--Ricci curvature} of an edge $(u,v) \in E$ is
\[
\kappa(u,v) = 1 - \frac{W_1(\mu_u, \mu_v)}{d(u,v)},
\]
where $W_1$ denotes the Wasserstein-1 (earth mover's) distance.
\end{definition}

Positive curvature indicates convergent geometry (neighborhoods of $u$ and $v$ are close); negative curvature indicates divergent geometry characteristic of bridges, bottlenecks, or saddle-like structure between distinct communities.

\paragraph{Ricci flow on graphs.}
Discrete Ricci flow iteratively updates edge weights \cite{ni2019community}:
\[
w_{uv}^{(t+1)} = w_{uv}^{(t)} - \eta \cdot \kappa^{(t)}(u,v) \cdot w_{uv}^{(t)},
\]
where $\eta > 0$ is a step size.
After $T$ iterations, shortest-path distances on the reweighted graph define the Ricci backend metric $D^{(\text{ricci})}$.

\paragraph{Why Ricci may help on bottleneck and shortcut structures.}
When a $k$NN graph contains shortcut edges between biologically distinct subpopulations---e.g., between LUAD subtypes that are separated by a thin manifold bridge in expression space---these edges typically have negative Ollivier--Ricci curvature.
Ricci flow upweights negatively curved (bridge) edges, effectively increasing the filtration threshold at which they create topological connections.
This delays the merger of distinct components, preserving genuine cycles at the expense of noise-induced shortcuts.
The net effect is that the persistence barcode under Ricci geometry may have longer-lived, more biologically meaningful $H_1$ bars---though this effect is data-dependent and not guaranteed.

\begin{remark}[Non-monotonicity of Ricci benefits]\label{rem:ricci_nonmono}
Ricci flow does not uniformly improve downstream performance.
It helps when curvature correction suppresses misleading shortcut edges (as on LUAD).
It does not help when Euclidean neighborhoods already capture the relevant geometry (as on GSE161711's smaller, more homogeneous cohort), or when the dataset is too small for the geometry effect to dominate split variance.
This non-monotonicity motivates treating geometry as a diagnostic axis.
\end{remark}

\emph{Implementation status: fully implemented and benchmarked.}

\subsection{Diffusion geometry, heat kernels, and manifold learning}
\label{sec:diffusion}

Diffusion geometry replaces pairwise distances with quantities derived from a graph Laplacian or heat kernel, capturing connectivity that persists across random walks of varying length \cite{coifman2006diffusion}.

\begin{definition}[Diffusion distance]\label{def:diffusion}
Let $W$ be a weight matrix derived from a Gaussian kernel $W_{ij} = \exp(-\|x_i - x_j\|^2 / \sigma^2)$, and let $P = D_W^{-1} W$ be the row-normalized Markov transition matrix where $D_W = \operatorname{diag}(\sum_j W_{ij})$.
The \emph{diffusion distance} at time $t$ is
\[
D^{(\text{diff})}_{t}(i,j) = \Bigl( \sum_{\ell \ge 1} e^{-2t\mu_\ell} (\psi_\ell(i) - \psi_\ell(j))^2 \Bigr)^{1/2},
\]
where $\mu_\ell$ and $\psi_\ell$ are eigenvalues and right eigenvectors of $P$.
\end{definition}

\begin{remark}[Multiscale property and omics relevance]\label{rem:diffusion_multiscale}
The diffusion time parameter $t$ controls the scale of geometry: small $t$ recovers local Euclidean-like distances; large $t$ emphasizes global connectivity.
In omics, where manifold-like structure is common (differentiation trajectories, cell-state continua \cite{haghverdi2016}), diffusion distances are often more biologically appropriate than Euclidean distances because they respect the intrinsic geometry of the data manifold rather than ambient-space proximity.
Diffusion distances can equivalently be expressed through the heat kernel $H_t = e^{-tL}$ as $D^{(\text{diff})}_t(i,j)^2 = H_t(i,i) + H_t(j,j) - 2H_t(i,j)$, connecting persistent homology to graph spectral methods.
\end{remark}

\paragraph{Why diffusion geometry may better denoise local noise.}
In a persistent-homology context, diffusion backends produce filtrations that are selectively sensitive to structural scales controlled by $t$.
Local noise (e.g., measurement error, dropout in scRNA-seq) creates high-frequency perturbations that the diffusion operator damps at even moderate $t$, while preserving the low-frequency signal that drives genuine topological features.
This makes diffusion geometry a natural denoising backend for noisy omics data.

\emph{Implementation status: mathematically developed; not experimentally instantiated.}

\subsection{Information geometry and PHATE-like distances}
\label{sec:phate}

PHATE (Potential of Heat-diffusion for Affinity-based Trajectory Embedding) \cite{moon2019phate} transforms diffusion operators into potential coordinates via a log transformation.

\begin{definition}[PHATE-like potential distance]\label{def:phate}
Let $P^t$ be the $t$-step transition matrix.
The \emph{potential representation} of sample $i$ is $U_i = -\log(P^t_{i,\cdot} + \epsilon)$, where $\epsilon > 0$ is a regularization constant.
The \emph{potential distance} is
\[
D^{(\text{phate})}(i,j) = \|U_i - U_j\|_2.
\]
\end{definition}

\begin{remark}[Information-geometric interpretation]\label{rem:phate_info}
The log transformation converts transition probabilities into surprisal (information content).
The potential distance therefore measures how differently the random walks emanating from $i$ and $j$ explore the graph, related to the symmetrized KL divergence between diffusion distributions.
In information-geometric terms, $D^{(\text{phate})}$ approximates a Fisher--Rao-like distance on the statistical manifold of diffusion distributions.
\end{remark}

\paragraph{Why PHATE-like distances may align with developmental or cellular trajectories.}
The potential distance preserves both local neighborhood structure and coarse global organization, making it suitable for data with branching trajectories or state transitions where loops may arise from biological cycles (e.g., cell cycles, metabolic oscillations).
In single-cell omics, where cells are often sampled along continuous differentiation paths, PHATE-like distances can produce filtrations that better respect the trajectory geometry than Euclidean distances, potentially revealing topological features (loops at branch points, voids between lineages) that correspond to genuine biological programs.

\emph{Implementation status: mathematically developed; not experimentally instantiated.}

\subsection{Robust topology via distance-to-measure (DTM)}
\label{sec:dtm}

The distance-to-measure (DTM) framework \cite{chazal2011dtm, chazal2017robust} replaces pairwise distances with density-smoothed variants that are robust to outliers and density heterogeneity.

\begin{definition}[Distance to measure]\label{def:dtm}
For a point cloud $X$ with empirical measure $\mu_X = \frac{1}{n}\sum_{i=1}^n \delta_{x_i}$ and a mass parameter $m_0 \in (0,1]$, the \emph{distance-to-measure function} at point $x$ is
\[
d_{\mu_X, m_0}(x) = \Bigl( \frac{1}{\lceil m_0 n \rceil} \sum_{j=1}^{\lceil m_0 n \rceil} \|x - x_{(j)}\|^2 \Bigr)^{1/2},
\]
where $x_{(1)}, \ldots, x_{(n)}$ are the points of $X$ sorted by distance to $x$.
\end{definition}

\begin{proposition}[Stability of DTM \cite{chazal2011dtm}]\label{prop:dtm}
If $\mu$ and $\nu$ are two probability measures on $\mathbb{R}^p$, then
\[
\|d_{\mu, m_0} - d_{\nu, m_0}\|_\infty \le \frac{1}{\sqrt{m_0}} W_2(\mu, \nu),
\]
where $W_2$ is the Wasserstein-2 distance.
\end{proposition}

This stability guarantee means DTM-based filtrations are robust to perturbations in the data-generating distribution, including the addition or removal of outlier samples---particularly relevant for omics data with batch effects, dropout noise, or technical outliers.

\paragraph{DTM-filtrations and robustness to density irregularity.}
A DTM-weighted distance $D^{(\text{dtm})}_{ij} = \|x_i - x_j\|^2 + d_{\mu_X, m_0}(x_i)^2 + d_{\mu_X, m_0}(x_j)^2$ delays the inclusion of outlier-adjacent simplices, suppressing short-lived topological noise from density spikes or isolated points.
In omics, where outlier samples can arise from technical artifacts (failed library preparations, doublets in scRNA-seq), DTM-based filtrations can produce more biologically meaningful barcodes by preventing these artifacts from dominating the persistence landscape.

\emph{Implementation status: mathematically developed; not experimentally instantiated.}

% ============================================================
\section{Experimental Design Under a Resource-Bounded Policy}
\label{sec:experimental}
% ============================================================

\subsection{Datasets and roles}

We use four datasets with distinct roles in the evaluation:
\begin{itemize}[leftmargin=1.5em]
\item \textbf{TCGA-LUAD} (lung adenocarcinoma, $n \approx 500$, $p \approx 20{,}000$): the main benchmark anchor with multiseed paired statistics across 18 random train/test splits, providing sufficient power for paired $t$-tests and Wilcoxon signed-rank tests.
\item \textbf{GSE161711} (CLL/venetoclax bulk RNA-seq, $n = 57$): disease-aligned bulk validation with 5-fold cross-validation and full six-axis ablation sweep (geometry backend, Ricci iterations, $k$, $\lambda$, normalization, top-$K$).
\item \textbf{BRCA} (breast cancer, $n \approx 1{,}000$, $p \approx 20{,}000$): supplementary supporting evidence with 2-fold cross-validation.
\item \textbf{GSE165087} (CLL/RS-like scRNA-seq, $n_{\text{bounded}} \in \{320, 400\}$): bounded single-cell heterogeneity validation under strict resource constraints after three larger configurations failed.
\end{itemize}

GSE161711 and GSE165087 are public GEO substitutes for originally intended controlled-access datasets that were unavailable in the present execution context.
They are scientifically reasonable substitutes (CLL disease setting, drug-resistance biology) but not identical to the original targets.

\subsection{Execution envelope and resource policy}

All experiments were conducted on an HPC cluster (Hellbender) under explicit resource constraints.
LUAD and GSE161711 completed within modest allocations ($\le$48\,GB, $\le$40\,min).
For single-cell data, three progressively smaller configurations exceeded available memory or walltime before bounded configurations succeeded.
This failure sequence defines the \emph{execution envelope} of the paper: the set of experiments that were attempted, the resources they consumed, and the outcomes.
Complete Slurm job histories, peak RSS observations, and runtime details are documented in the Supplement.

\paragraph{Resource policy.}
Rather than silently omitting failed runs, this paper adopts an explicit resource-bounded policy: completed outputs are reused in place, heavy failed runs are not naively repeated, and new single-cell passes are only permitted when they are strictly smaller than every prior failure.
This policy shapes both the evidence presented and the limitations stated.

\subsection{Evaluation protocol}

Each dataset follows the same pipeline: geometry backend $\to$ Rips filtration $\to$ persistent homology $\to$ sparse representative $\to$ TIP $\to$ feature selection $\to$ weighted-F1 classification.
Dataset-appropriate replication:
\begin{itemize}[leftmargin=1.5em]
\item LUAD: 18 random splits (paired multiseed statistics).
\item GSE161711: 5-fold CV.
\item BRCA: 2-fold CV.
\item GSE165087: 1 bounded split, 1 bootstrap.
\end{itemize}

\subsection{Completed versus bounded versus conceptual experiments}

A central integrity constraint of this paper is the three-way separation:
\begin{enumerate}[leftmargin=1.5em]
\item \textbf{Completed core experiments}: LUAD (full benchmark + multiseed + ablation), GSE161711 (5-fold CV + full ablation), BRCA (supplementary 2-fold).
\item \textbf{Bounded sensitivity analyses}: GSE165087 manuscript-safe ($n_{\text{cells}} = 320$) and ablation-safe ($n_{\text{cells}} = 400$, lambda-only sensitivity).
\item \textbf{Future backend completions}: diffusion, PHATE-like, DTM implementation and benchmarking under larger resources.
\end{enumerate}
No empirical claim in this paper depends on an unimplemented backend or an incomplete experiment.

% ============================================================
\section{Results}
\label{sec:results}
% ============================================================

\subsection{Does sparsity help representative interpretability?}

The primary evidence that sparsity improves interpretability comes from two sources: the modest but consistent CC-SPR advantage over harmonic PH on datasets where the topological signal is resolvable, and the TIP diagnostics that show how sparse representatives concentrate feature influence.

On GSE161711, CC-SPR achieves mean F1 $= 0.603$, exceeding harmonic PH ($0.585$, $\Delta = +0.018$).
On the ablation-safe GSE165087 configuration, CC-SPR achieves F1 $= 0.823$ versus harmonic $= 0.798$ ($\Delta = +0.025$).
These gains are modest but directionally consistent: sparse representatives project to a more compact gene set, reducing the dilution effect of including many low-signal genes in the feature vector.

The TIP diagnostics (Figures~\ref{fig:tip_luad} and \ref{fig:tip_bulk}) provide complementary evidence.
When sparse representatives and TIP concentration agree---a few features repeatedly selected across bootstrap replicates---the method yields a stronger feature-level interpretation than a dense representative alone.

\begin{table}[h]
\centering
\small
\begin{tabular}{llccc}
\toprule
Dataset & Model & Mean F1 & Std & $N$ \\
\midrule
LUAD & standard & 0.7670 & 0.1386 & 2 \\
LUAD & harmonic & 0.4044 & 0.0881 & 2 \\
LUAD & ccspr (Ricci) & 0.3133 & 0.1603 & 2 \\
\midrule
GSE161711 & standard & 0.9711 & 0.0404 & 5 \\
GSE161711 & harmonic & 0.5850 & 0.0307 & 5 \\
GSE161711 & ccspr (Ricci) & 0.6028 & 0.1562 & 5 \\
\midrule
GSE165087 (bounded) & standard & 0.9254 & --- & 1 \\
GSE165087 (bounded) & harmonic & 0.9254 & --- & 1 \\
GSE165087 (bounded) & ccspr (Ricci) & 0.9254 & --- & 1 \\
\midrule
GSE165087 (ablation-safe) & standard & 0.888 & --- & 1 \\
GSE165087 (ablation-safe) & harmonic & 0.798 & --- & 1 \\
GSE165087 (ablation-safe) & ccspr (Ricci) & 0.823 & --- & 1 \\
\midrule
BRCA & standard & 0.6112 & 0.0073 & 2 \\
BRCA & harmonic & 0.4788 & 0.0891 & 2 \\
BRCA & ccspr (Ricci) & 0.4096 & 0.0142 & 2 \\
\bottomrule
\end{tabular}
\caption{Main benchmark results across all datasets.
Standard baselines dominate on raw F1.
CC-SPR exceeds harmonic PH on GSE161711 ($+0.018$) and ablation-safe GSE165087 ($+0.025$).
On LUAD and BRCA, all topology-derived methods trail the standard baseline.
The bounded GSE165087 configuration produces a three-way tie due to the small cell budget.}
\label{tab:mainbench}
\end{table}

\paragraph{Interpretation.}
The standard baseline's dominance reflects a well-known pattern: when class structure aligns with variance-dominant axes (as in TCGA expression subtypes), linear classifiers on high-variance genes are hard to beat.
Topology-derived features capture complementary structural information---loop-like patterns in sample space---that may correspond to regulatory circuits, microenvironment heterogeneity, or subtype-transition geometry.
The value of CC-SPR is not predictive dominance but rather that sparse representatives and TIP profiles identify a compact, stable gene subset associated with the dominant $H_1$ cycle, providing a topology-grounded alternative to purely variance-based feature selection.

\subsection{LUAD multiseed paired statistics}

The multiseed analysis provides more informative evidence than the two-split summary because it measures paired differences across 18 repeated splits, controlling for split-specific data variation.

\begin{table}[h]
\centering
\small
\begin{tabular}{lcccc}
\toprule
Comparison & $n$ & Mean $\Delta$ & Paired $t$ ($p$) & Wilcoxon ($p$) \\
\midrule
CC-SPR vs harmonic & 18 & $+0.027$ & 0.514 & 0.495 \\
CC-SPR vs standard & 18 & $-0.186$ & $2.7 \times 10^{-4}$ & $3.3 \times 10^{-4}$ \\
Ricci sparse vs Eu.\ sparse & 18 & $+0.046$ & 0.308 & 0.265 \\
\bottomrule
\end{tabular}
\caption{LUAD multiseed paired statistics ($n=18$ splits).
CC-SPR significantly underperforms the standard baseline ($p < 0.001$).
The Ricci-vs-Euclidean sparse comparison shows a positive mean delta ($+0.046$) but is not significant at this sample size ($p = 0.308$).}
\label{tab:paired}
\end{table}

\paragraph{The Ricci-vs-Euclidean sparse result.}
The $+0.046$ mean delta in favor of Ricci sparse is substantively interesting but statistically nonsignificant.
On average, Ricci curvature correction modestly improves sparse representative quality on LUAD, but the effect is smaller than split-to-split variability.
This is exactly the backend-sensitive behavior the framework is designed to expose: geometry matters, but its effect size is dataset-dependent and must be diagnosed empirically rather than assumed.

\begin{figure}[h]
\centering
\includegraphics[width=0.49\linewidth]{../results/figures/f1_bars_ci.png}
\includegraphics[width=0.49\linewidth]{../results/figures/paired_deltas_luad_multiseed.png}
\caption{\textbf{Left}: LUAD benchmark showing the predictive gap between the standard baseline and topology-derived models.
\textbf{Right}: Paired multiseed deltas across 18 splits.
The CC-SPR vs standard comparison is significantly negative; the Ricci vs Euclidean sparse comparison is positive but not significant.
Paired statistics are more informative than summary means because they control for split-specific data variation.}
\label{fig:luad}
\end{figure}

\subsection{When does geometry help? GSE161711 as disease-aligned validation}

On GSE161711 (CLL patients treated with venetoclax, $n = 57$), CC-SPR achieves mean F1 $= 0.603$, modestly exceeding harmonic PH ($0.585$, $\Delta = +0.018$) while trailing the standard baseline ($0.971$).
This dataset is biologically closer to the motivating CLL setting than LUAD: venetoclax resistance involves rewiring of apoptotic pathways (BCL-2 family interactions, metabolic reprogramming) that may introduce loop-like structure in expression space as cells transition between sensitive and resistant states.

The modest CC-SPR advantage over harmonic PH suggests that sparse representatives can capture some of this transitional geometry, though the effect is small relative to the dominant variance signal exploited by the standard classifier.
The TIP diagnostics (Figure~\ref{fig:tip_bulk}) reveal that Euclidean and Ricci backends concentrate feature influence on partially overlapping but distinct gene subsets, indicating that geometry choice genuinely affects which biological signals the topological pipeline highlights.

\begin{figure}[h]
\centering
\includegraphics[width=0.49\linewidth]{../results/figures/tip_eu_vs_ricci_cll_venetoclax_gse161711.png}
\includegraphics[width=0.49\linewidth]{../results/figures/h1_lifetime_eu_vs_ricci_cll_venetoclax_gse161711.png}
\caption{GSE161711 topological diagnostics.
\textbf{Left}: TIP comparison between Euclidean and Ricci backends shows geometry affects feature concentration and which genes are selected.
\textbf{Right}: $H_1$ lifetime comparison shows the dominant topological signal changes under different geometry assumptions.}
\label{fig:tip_bulk}
\end{figure}

\subsection{When does Ricci help, and when does it not?}
\label{sec:ricci_analysis}

The completed empirical evidence reveals three distinct regimes:

\begin{enumerate}[leftmargin=1.5em]
\item \textbf{LUAD}: Ricci sparse achieves a positive mean delta ($+0.046$) over Euclidean sparse in paired multiseed analysis, but the effect is not statistically significant ($p = 0.308$).
The ablation sweep shows that the Ricci benefit is concentrated at higher iteration counts ($T=10$ preferred) and is sensitive to $\lambda$.

\item \textbf{GSE161711}: The ablation reveals that Euclidean geometry is preferred over Ricci (Eu: $0.600$ vs Ricci: $0.532$), and zero Ricci iterations outperform positive iterations.
The smaller cohort ($n=57$) likely does not have the bridge/shortcut structure that Ricci flow is designed to correct.

\item \textbf{GSE165087 (bounded)}: All methods tie at F1 $= 0.925$ on the bounded split.
The ablation-safe configuration ($n_{\text{cells}}=400$) shows CC-SPR $> $ harmonic, but does not resolve Euclidean vs Ricci.
\end{enumerate}

The underlying mechanism is interpretable: Ricci flow reweights edges by curvature, contracting positively curved regions and expanding negatively curved ones.
When the data contains shortcut edges that create spurious topological features (as may occur in LUAD's heterogeneous subtype landscape), Ricci correction can sharpen the true signal.
When the dominant topology is already well-captured by Euclidean neighborhoods (as in GSE161711's smaller, more homogeneous cohort), the additional geometric processing adds noise without benefit.

\begin{figure}[h]
\centering
\includegraphics[width=0.72\linewidth]{../results/figures/geometry_backend_summary.png}
\caption{Implemented geometry comparison across datasets.
Euclidean harmonic vs Ricci sparse on the three datasets with completed outputs.
Diffusion, PHATE-like, and DTM backends are not included because they are not experimentally instantiated.}
\label{fig:geombar}
\end{figure}

\subsection{What do the topological diagnostics reveal?}

\paragraph{Lambda sensitivity confirms sparsity is a real modeling axis.}
The LUAD $\lambda$ sweep (Figure~\ref{fig:diagnostics}, left) shows that predictive performance varies meaningfully across sparsity settings, confirming that the elastic-net penalty in \eqref{eq:ccspr} is not decorative.
Different $\lambda$ values produce different representative support sizes and different downstream classifier performance.
This validates the methodological claim that sparsity should be a tunable parameter in topological feature attribution.

\paragraph{$H_1$ lifetime and prominence reveal geometry's effect on topology.}
The $H_1$ lifetime and prominence diagnostics (Figures~\ref{fig:diagnostics}, \ref{fig:tip_luad}) show that geometry changes the persistence landscape itself, not only the downstream classifier.
On GSE161711, the prominence of the dominant $H_1$ bar changes under Ricci geometry, indicating that curvature correction modifies which topological features are most salient.
These diagnostics motivate the method even when classifier rankings do not settle the question by themselves.

\begin{figure}[h]
\centering
\includegraphics[width=0.49\linewidth]{../results/figures/lambda_robustness.png}
\includegraphics[width=0.49\linewidth]{../results/figures/h1_prominence_eu_vs_ricci_cll_venetoclax_gse161711.png}
\caption{\textbf{Left}: LUAD $\lambda$ sensitivity. Representative sparsity is a real modeling axis---different $\lambda$ values produce different predictive performance.
\textbf{Right}: GSE161711 $H_1$ prominence. Geometry changes the salience of the dominant topological signal, even when classifier rankings remain conservative.}
\label{fig:diagnostics}
\end{figure}

\begin{figure}[h]
\centering
\includegraphics[width=0.49\linewidth]{../results/figures/tip_eu_vs_ricci.png}
\includegraphics[width=0.49\linewidth]{../results/figures/h1_lifetime_eu_vs_ricci.png}
\caption{LUAD TIP and $H_1$ lifetime diagnostics comparing Euclidean and Ricci backends.
TIP concentration differences reveal how geometry affects feature-level interpretability; lifetime differences show how geometry modifies the persistence landscape.}
\label{fig:tip_luad}
\end{figure}

\subsection{Ablation analysis: LUAD and GSE161711}

The full ablation story is carried by LUAD and GSE161711, which have completed cached sweeps over six hyperparameter axes.

\begin{table}[h]
\centering
\small
\begin{tabular}{llccc}
\toprule
Dataset & Ablation axis & Best value & Mean F1 & $n$ \\
\midrule
LUAD & distance mode & ricci & 0.463 & 2 \\
LUAD & Ricci iterations & 10 & 0.551 & 2 \\
LUAD & $k$ & 10 & 0.405 & 2 \\
LUAD & $\lambda$ & $10^{-7}$ & 0.594 & 2 \\
LUAD & normalize & std & 0.522 & 2 \\
LUAD & top-$K$ & 10 & 0.667 & 2 \\
\midrule
GSE161711 & distance mode & euclidean & 0.600 & 5 \\
GSE161711 & Ricci iterations & 0 & 0.578 & 5 \\
GSE161711 & $k$ & 5 & 0.632 & 5 \\
GSE161711 & $\lambda$ & $10^{-4}$ & 0.590 & 5 \\
GSE161711 & normalize & std & 0.596 & 5 \\
GSE161711 & top-$K$ & 50 & 0.806 & 5 \\
\bottomrule
\end{tabular}
\caption{Best completed ablation settings from cached sweeps.
LUAD prefers Ricci geometry, higher iteration count, compact top-$K$, and moderate $\lambda$.
GSE161711 prefers Euclidean geometry, zero iterations, broad top-$K$, and higher $\lambda$.
This cross-dataset divergence is the backend-sensitive behavior the paper claims.}
\label{tab:ablationbest}
\end{table}

\begin{figure}[h]
\centering
\includegraphics[width=0.78\linewidth]{../results/figures/ablation_profiles_luad_gse161711.png}
\caption{Cross-dataset ablation profiles from completed LUAD and GSE161711 sweeps.
The pipeline exhibits dataset-specific sensitivity to every hyperparameter axis, supporting the claim that geometry and sparsity are genuine modeling decisions.}
\label{fig:ablationprofiles}
\end{figure}

\paragraph{Cross-dataset divergence as evidence for backend sensitivity.}
The most informative finding from the ablation analysis is that LUAD and GSE161711 prefer \emph{different} settings across every axis.
LUAD benefits from Ricci geometry, more iterations, and compact top-$K$; GSE161711 benefits from Euclidean geometry, zero iterations, and broad top-$K$.
This divergence is not a failure---it is the core evidence that geometry and sparsity parameters are genuine modeling decisions whose optimal values depend on dataset structure.
A method that always preferred the same settings would be insensitive to geometry, undermining the scientific claim.

\subsection{What is gained biologically even when prediction does not dominate?}

Even when standard baselines win on F1, the topological pipeline contributes if it yields a stable sparse explanatory object.
On LUAD, the sparse representatives and TIP profiles identify a compact gene subset associated with the dominant $H_1$ cycle.
This subset is interpretable as a topology-grounded feature selection: genes that repeatedly appear in TIP across bootstraps are those whose coordinated variation defines a persistent loop in sample space.

On GSE161711, the CC-SPR advantage over harmonic PH ($+0.018$) is accompanied by TIP diagnostics showing that Euclidean and Ricci backends highlight partially overlapping but distinct gene sets.
This means geometry choice is not merely a technical parameter---it determines which biological signals the topological pipeline surfaces.
A practitioner can compare TIP profiles across backends to identify genes that are robustly associated with topology regardless of geometry (core signal) versus genes that are geometry-dependent (backend-specific signal).

% ============================================================
\section{Resource-Bounded Single-Cell Validation}
\label{sec:scrna}
% ============================================================

The single-cell story must be handled with care and full transparency.

\subsection{Execution history}

We successfully produced the public substitute dataset and preprocessing artifact \texttt{data/cll\_rs\_scrna/data.h5ad} (GSE165087), but larger downstream experiment configurations repeatedly failed under cluster limits.

\begin{table}[h]
\centering
\small
\begin{tabular}{lllll}
\toprule
Run type & State & Elapsed & Requested memory & Peak RSS \\
\midrule
scRNA full & OUT\_OF\_MEMORY & 16:43:25 & 96\,GB & 100.3\,GB \\
scRNA fast-track & TIMEOUT & 20:00:06 & 96\,GB & 75.2\,GB \\
scRNA low-memory & OUT\_OF\_MEMORY & 11:25:59 & 48\,GB & 49.7\,GB \\
scRNA manuscript-safe & COMPLETED & 00:07:49 & 32\,GB & 2.9\,GB \\
scRNA ablation-safe & COMPLETED & --- & 32\,GB & ${\sim}$9.3\,GB \\
\bottomrule
\end{tabular}
\caption{Execution envelope for single-cell experiments.
Three configurations failed, establishing the resource boundary.
Two bounded configurations completed, providing feasibility evidence and limited sensitivity analysis.}
\label{tab:scRNAresource}
\end{table}

\paragraph{Failure analysis.}
The full scRNA run OOMed at ${\sim}$100.3\,GB RSS under a 96\,GB request after 16+ hours.
A reduced fast-track configuration timed out after 20 hours.
A further reduced low-memory configuration still OOMed at ${\sim}$49.7\,GB under a 48\,GB request.
The memory bottleneck arises from the combination of $k$NN graph construction, Ricci flow iterations, persistent homology computation, and bootstrap loops on the full scRNA dataset---a chain of $O(n^2)$ or worse operations on $n \gg 10^3$ cells.

\paragraph{Code-level correction.}
A key code-level correction enabled the bounded runs: the preprocessed single-cell loader originally reused the full stored PCA matrix (\texttt{X\_pca} in the \texttt{.h5ad} file) regardless of the \texttt{n\_pcs} parameter.
After correction, the loader truncates to the requested number of principal components, ensuring that the manuscript-safe run genuinely operates in a 20-dimensional space rather than the 50-dimensional space stored in the preprocessed file.

\subsection{Bounded results}

\paragraph{Manuscript-safe configuration} ($n_{\text{cells}} = 320$, $n_{\text{pcs}} = 20$, 1 split, 1 bootstrap):
All three models (standard, harmonic, CC-SPR) achieve identical weighted F1 $= 0.9254$.
This is expected for a small bounded split where the dominant signal is the class structure itself; the configuration is too small to resolve method differences.

\paragraph{Ablation-safe configuration} ($n_{\text{cells}} = 400$, $n_{\text{pcs}} = 25$, 1 split, 1 bootstrap):
Standard F1 $= 0.888$, harmonic $= 0.798$, CC-SPR $= 0.823$ ($+0.025$ vs harmonic).
Lambda sensitivity: $\lambda = 10^{-6}$ yields F1 $= 0.847$; $\lambda = 10^{-5}$ yields F1 $= 0.850$.
The ablation-safe configuration provides more informative results than the manuscript-safe tie, showing CC-SPR exceeding harmonic while trailing standard---consistent with the bulk datasets.

\subsection{Interpretation}

These bounded results establish feasibility and directional consistency with the bulk datasets, but they do not substitute for comprehensive single-cell benchmarking.
This paper is an analytical methods study, and the single-cell component serves as a bounded heterogeneity probe rather than a primary benchmark axis.
Larger single-cell evaluations---with more cells, more splits, more bootstraps, and full ablation grids---remain future work contingent on memory optimization or access to larger-memory nodes.

% ============================================================
\section{Biological Interpretation}
\label{sec:biology}
% ============================================================

A persistent-homology method for omics is only useful if its outputs connect to biology.
This section interprets the CC-SPR results in biological terms, using existing outputs and disciplined inference.

\subsection{LUAD subtype structure and topological signal}

TCGA-LUAD contains well-characterized molecular subtypes (terminal respiratory unit / TRU, proximal-inflammatory / PI, proximal-proliferative / PP) that define the dominant variance structure in expression space.
The standard baseline exploits this variance structure directly, which is why it achieves the highest F1.

The CC-SPR pipeline detects $H_1$ cycles in sample space that correspond to loop-like relationships among these subtypes.
Such loops can arise when subtype boundaries are not sharp---when transitional samples create bridges between clusters---or when a regulatory program introduces periodic-like variation (e.g., cell-cycle-associated expression in proliferative subtypes).
The sparse representatives project this topological signal to a compact gene set.
Although we do not have cached pathway enrichment results, the fact that TIP profiles show concentrated feature influence suggests that the selected genes are not random but reflect a coherent molecular signal.
Connecting TIP-selected genes to pathway databases (KEGG, Reactome, MSigDB) is a natural and important next step that we identify as future work.

\subsection{GSE161711: disease-aligned bulk structure}

GSE161711 profiles CLL patients treated with venetoclax, a BCL-2 inhibitor.
Resistance to venetoclax involves complex rewiring of apoptotic and metabolic programs: cells may upregulate alternative anti-apoptotic proteins (MCL-1, BCL-XL), shift metabolic dependencies, or activate pro-survival signaling (NF-$\kappa$B, PI3K/AKT).
This rewiring can create continuous transitions in expression space that manifest as topological features---loops or bridges connecting sensitive and resistant states.

The CC-SPR advantage over harmonic PH on this dataset ($+0.018$) is modest but biologically plausible: sparse representatives may be better at isolating the gene module associated with this resistance-transition topology, while dense representatives dilute the signal across many genes.
The TIP comparison between Euclidean and Ricci backends shows partially overlapping but distinct gene selections, suggesting that geometry choice determines which aspect of the resistance biology the topological pipeline highlights.

\subsection{Bounded GSE165087: cellular heterogeneity interpretation}

GSE165087 profiles CLL and Richter syndrome cells, which exhibit substantial intra-tumor heterogeneity.
Even in the bounded configuration ($n_{\text{cells}} = 400$), the CC-SPR advantage over harmonic PH ($+0.025$) suggests that sparse representatives can capture heterogeneity-related structure more effectively than dense ones.

In single-cell data, $H_1$ cycles may correspond to cell-state transitions (e.g., cycling vs quiescent CLL cells), regulatory feedback loops within a subpopulation, or sampling-induced loops from uneven coverage of a continuous manifold.
The bounded nature of our evaluation means we cannot distinguish these possibilities, but the directional consistency with bulk results (CC-SPR $>$ harmonic) provides a plausible signal.
Larger single-cell studies with pathway enrichment and pseudotime analysis would be needed to resolve the biological interpretation.

\subsection{Why topology-aware sparse representatives identify concentrated disease-relevant modules}

The fundamental argument is structural: a persistent $H_1$ cycle in sample space corresponds to a set of samples whose pairwise distances create a loop in the Rips filtration.
When this cycle is projected to feature space through the data matrix, the resulting importance weights identify genes whose coordinated variation ``traces'' the loop.
Dense representatives spread this signal across many genes; sparse representatives concentrate it on a few.

In disease biology, concentrated gene modules are more interpretable because they suggest specific pathways or regulatory programs rather than diffuse transcriptome-wide effects.
A topological signal that localizes to 15 genes in an apoptotic pathway is more biologically actionable than one that implicates 200 genes at similar low weights.
CC-SPR's elastic-net objective directly promotes this concentration, making the method better suited to biological hypothesis generation than dense harmonic representatives.

\subsection{How CC-SPR improves interpretability even when predictive gains are modest}

The interpretability value of CC-SPR is independent of its predictive performance.
Even on LUAD, where CC-SPR trails the standard baseline by a significant margin, the sparse representatives and TIP profiles provide information that the standard classifier does not: they identify which genes are repeatedly associated with the dominant topological feature under resampling, and they show how this association changes under different geometry assumptions.

This diagnostic value is analogous to how unsupervised clustering can provide biological insights even when it does not improve supervised classification: both capture structure that is orthogonal to the class label.
The practical recommendation is to use CC-SPR not as a replacement for standard classifiers but as a complementary analysis that provides topology-grounded feature attribution and geometry sensitivity diagnostics.

% ============================================================
\section{Limitations and Future Directions}
\label{sec:limitations}
% ============================================================

This paper has six practical limitations that are part of the contribution because they distinguish a reproducible resource-aware method paper from a benchmark narrative.

\begin{enumerate}[leftmargin=1.5em]
\item \textbf{Standard baselines win on prediction.}
On LUAD, BRCA, and GSE161711, non-topological baselines achieve the highest weighted F1.
CC-SPR's contribution is methodological (interpretable sparse representatives, geometry control, stability diagnostics), not predictive dominance.
\emph{Future work}: integrating topology-derived features as additional inputs to conventional classifiers, rather than using them as standalone predictors.

\item \textbf{The backend family exceeds the implementation.}
Euclidean and Ricci are implemented and benchmarked.
Diffusion, PHATE-like, and DTM are mathematically developed (Sections~\ref{sec:diffusion}--\ref{sec:dtm}) and compared conceptually (Table~\ref{tab:hphcompare}), but not experimentally instantiated.
No empirical claim depends on an unimplemented backend.
\emph{Future work}: implementing diffusion, PHATE, and DTM backends; benchmarking on the same datasets; studying how each geometry affects TIP concentration and biological interpretability.

\item \textbf{Ricci benefit is non-monotone and dataset-dependent.}
LUAD shows a modest positive Ricci effect; GSE161711 shows a negative one.
This is the expected behavior of a backend-sensitive framework, not a failure.
\emph{Future work}: larger paired-comparison studies to determine under what data-structural conditions Ricci helps.

\item \textbf{Single-cell evidence is resource-bounded.}
Three scRNA configurations failed under reasonable cluster allocations.
The manuscript-safe and ablation-safe runs are sufficient for feasibility validation but not for comprehensive benchmarking.
\emph{Future work}: memory-optimized implementations (sparse distance matrices, approximate persistent homology, chunked Ricci flow); evaluation on larger cell budgets with multi-split, multi-bootstrap designs.

\item \textbf{Public substitutes replace controlled-access datasets.}
GSE161711 and GSE165087 are scientifically reasonable but not identical to the originally intended data.
\emph{Future work}: evaluation on the original controlled-access datasets when available.

\item \textbf{Representative diagnostics are partial.}
Per-bootstrap support entropy and Gini coefficients would require rerunning topology loops; only TIP, $H_1$ lifetime, and prominence are available from cached outputs.
\emph{Future work}: caching per-bootstrap support vectors to enable post hoc entropy/support-size analysis.
\end{enumerate}

\paragraph{Broader future directions.}
Beyond the specific limitations above, several natural extensions remain:
\begin{itemize}[leftmargin=1.5em]
\item Connecting TIP-selected gene sets to pathway enrichment databases (KEGG, Reactome, MSigDB) to test whether topology-derived features align with known biology.
\item Survival stratification using TIP-derived features to assess clinical relevance.
\item Cross-tool harmonic PH comparison (e.g., maTilDA parity) to validate that the harmonic baseline is correctly implemented.
\item Larger bootstrap counts to study TIP convergence rates.
\item Extension to higher homology dimensions ($H_2$ and beyond) when dataset size permits.
\end{itemize}

% ============================================================
\section{Related Work}
\label{sec:related}
% ============================================================

\paragraph{Persistent homology in biology.}
Topological data analysis has been applied to gene expression \cite{rizvi2017}, protein structure \cite{xia2014}, tumor ecosystems \cite{vipond2021}, neuroscience \cite{giusti2015}, and evolutionary biology \cite{chan2013topology}.
Most applications use standard Euclidean persistent homology without representative cycle analysis.

\paragraph{Harmonic persistent homology.}
The harmonic PH framework \cite{harmonicPH} introduced the bridge from persistent classes to sample-level representatives, enabling feature attribution through projection.
CC-SPR builds directly on this work, preserving its scaffold while modifying the representative objective and generalizing the metric.

\paragraph{Ricci curvature on graphs.}
Ollivier--Ricci curvature \cite{ollivier2007, ollivier2009} and its discrete flow variant \cite{ni2019community} have been used for community detection, graph classification, and network analysis.
The connection to persistent homology through filtration perturbation (Proposition~\ref{prop:stability}) appears to be novel.

\paragraph{Diffusion and PHATE.}
Diffusion maps \cite{coifman2006diffusion} and PHATE \cite{moon2019phate} are widely used for dimensionality reduction in single-cell genomics.
Using diffusion or PHATE distances as persistent-homology backends is a natural extension that has been explored informally but, to our knowledge, not systematized within a sparse representative framework.

\paragraph{Distance to measure.}
DTM-based topology \cite{chazal2011dtm, chazal2017robust} provides robust persistent homology under noise.
Its integration into a representative-cycle framework with feature projection appears unexplored.

\paragraph{Sparse topology.}
Sparse representations in TDA have been studied through $\ell_1$ optimal cycles \cite{chen2011hardness, dey2011optimal}, but typically in a computational-geometry setting without the biological feature-projection layer that CC-SPR provides.

% ============================================================
\section{Conclusion}
\label{sec:conclusion}
% ============================================================

CC-SPR extends harmonic persistent homology in two directions: sparse elastic representatives that concentrate cycle support for interpretability, and a backend-aware framework that treats the sample-space metric as a selectable, stability-bounded parameter.
Together with topological influence profiles, these extensions produce a richer analysis object than dense Euclidean representatives alone---one that enables systematic comparison of how geometry affects topology-derived features.

The harmonic PH framework is the parent: every harmonic PH analysis is a special case of CC-SPR with $\lambda_1 = 0$ and a Euclidean backend.
CC-SPR's contribution is not to replace this parent but to extend it with explicit sparsity control, geometry-backend awareness, and bootstrap-aggregated stability diagnostics.

Our evaluation on four cancer genomics datasets confirms that:
\begin{itemize}[leftmargin=1.5em]
\item The framework is functional and produces consistent results across datasets.
\item Geometry sensitivity is real and dataset-dependent (LUAD prefers Ricci; GSE161711 prefers Euclidean).
\item Sparse representatives can modestly improve over dense harmonic ones (GSE161711: $+0.018$; bounded GSE165087: $+0.025$).
\item Standard non-topological baselines remain strongest on raw F1---a shared limitation of TDA-based methods that we state honestly.
\item The value of CC-SPR is methodological and diagnostic: compact gene-set identification, bootstrap stability assessment, and systematic geometry comparison.
\end{itemize}

The geometry-backend family is broader than the currently implemented code: Euclidean and Ricci are benchmarked; diffusion, PHATE-like, and DTM are mathematically developed and justified but not yet experimentally instantiated.
Single-cell experiments are resource-bounded.
Both represent natural next steps under larger compute budgets.
The completed evidence is sufficient to support the main paper's claims; the broader program remains a scientifically motivated roadmap.

\bibliographystyle{plain}
\begin{thebibliography}{25}

\bibitem{harmonicPH}
D.~Damiani, O.~Vipond, and H.~A.~Harrington.
Probing omics data via harmonic persistent homology.
\emph{bioRxiv}, 2024.
\texttt{doi:10.1101/2024.09.15.613148}.

\bibitem{edelsbrunner2010}
H.~Edelsbrunner and J.~Harer.
\emph{Computational Topology: An Introduction}.
AMS, 2010.

\bibitem{carlsson2009}
G.~Carlsson.
Topology and data.
\emph{Bulletin of the AMS}, 46(2):255--308, 2009.

\bibitem{rizvi2017}
A.~H.~Rizvi, P.~G.~Camara, E.~K.~Kandber, et al.
Single-cell topological RNA-seq analysis reveals insights into cellular differentiation and development.
\emph{Nature Biotechnology}, 35(6):551--560, 2017.

\bibitem{xia2014}
K.~Xia and G.-W.~Wei.
Persistent homology analysis of protein structure, flexibility, and folding.
\emph{International Journal for Numerical Methods in Biomedical Engineering}, 30(8):814--844, 2014.

\bibitem{vipond2021}
O.~Vipond, J.~A.~Bull, P.~S.~Sherland, et al.
Multiparameter persistent homology landscapes identify immune cell spatial patterns in tumors.
\emph{PNAS}, 118(41), 2021.

\bibitem{moon2019phate}
K.~R.~Moon, D.~van Dijk, Z.~Wang, et al.
Visualizing structure and transitions in high-dimensional biological data.
\emph{Nature Biotechnology}, 37(12):1482--1492, 2019.

\bibitem{zomorodian2005}
A.~Zomorodian and G.~Carlsson.
Computing persistent homology.
\emph{Discrete \& Computational Geometry}, 33(2):249--274, 2005.

\bibitem{friedman1998}
J.~Friedman.
Computing Betti numbers via combinatorial Laplacians.
\emph{Algorithmica}, 21(4):331--346, 1998.

\bibitem{zou2005elasticnet}
H.~Zou and T.~Hastie.
Regularization and variable selection via the elastic net.
\emph{JRSS Series B}, 67(2):301--320, 2005.

\bibitem{chazal2014persistence}
F.~Chazal, V.~de Silva, M.~Glisse, and S.~Oudot.
The structure and stability of persistence modules.
\emph{Springer}, 2016.

\bibitem{ollivier2007}
Y.~Ollivier.
Ricci curvature of metric spaces.
\emph{Comptes Rendus Mathematique}, 345(11):643--646, 2007.

\bibitem{ollivier2009}
Y.~Ollivier.
Ricci curvature of Markov chains on metric spaces.
\emph{Journal of Functional Analysis}, 256(3):810--864, 2009.

\bibitem{ni2019community}
C.-C.~Ni, Y.-Y.~Lin, F.~Luo, and J.~Gao.
Community detection on networks with Ricci flow.
\emph{Scientific Reports}, 9:9984, 2019.

\bibitem{coifman2006diffusion}
R.~R.~Coifman and S.~Lafon.
Diffusion maps.
\emph{Applied and Computational Harmonic Analysis}, 21(1):5--30, 2006.

\bibitem{haghverdi2016}
L.~Haghverdi, M.~B\"uttner, F.~A.~Wolf, F.~Buettner, and F.~J.~Theis.
Diffusion pseudotime robustly reconstructs lineage branching.
\emph{Nature Methods}, 13(10):845--848, 2016.

\bibitem{chazal2011dtm}
F.~Chazal, D.~Cohen-Steiner, and Q.~M\'erigot.
Geometric inference for probability measures.
\emph{Foundations of Computational Mathematics}, 11(6):733--751, 2011.

\bibitem{chazal2017robust}
F.~Chazal, B.~Fasy, F.~Lecci, B.~Michel, A.~Rinaldo, and L.~Wasserman.
Robust topological inference: distance to a measure and kernel distance.
\emph{JMLR}, 18(159):1--40, 2017.

\bibitem{gudhi2014}
The GUDHI Project.
GUDHI User and Reference Manual, 2014--2024.

\bibitem{giusti2015}
C.~Giusti, R.~Ghrist, and D.~S.~Bassett.
Two's company, three (or more) is a simplex.
\emph{Journal of Computational Neuroscience}, 41(1):1--14, 2016.

\bibitem{chan2013topology}
J.~M.~Chan, G.~Carlsson, and R.~Rabadan.
Topology of viral evolution.
\emph{PNAS}, 110(46):18566--18571, 2013.

\bibitem{chen2011hardness}
C.~Chen and D.~Freedman.
Hardness results for homology localization.
\emph{Discrete \& Computational Geometry}, 45(3):425--448, 2011.

\bibitem{dey2011optimal}
T.~K.~Dey, J.~Sun, and Y.~Wang.
Approximating loops in a shortest homology basis from point data.
\emph{SoCG}, 2010.

\end{thebibliography}

\end{document}
